To begin the discussion of our results, our numerical simulations reveal that for pristine graphene, the transmission coefficient generally remains very close to unity ($T \approx 0.99$), with a slight decrease as the barrier height increases. This slight sub-unity value arises from the multi-momentum composition of the Gaussian wave packet, which generates internal interference as distinct momentum components interact with the barrier simultaneously\cite{Staelens2021, MolgadoMex2018}. Conversely, the introduction of the Rashba spin-orbit interaction (SOIR) induces a more complex, non-monotonic variation in the transmission on the order of $10^{-2}$, highlighting the influence of spin-momentum coupling on transport through the barrier.

To understand these numerical trends observed in our time-dependent wave-packet simulations, we highlight that the expression for the current vector used in our analysis \eqref{eq:componentes} follows directly from the standard Dirac–Rashba Hamiltonian for graphene and has been reported in related contexts\cite{AvishaiPhysRevB2021}.
In this work, we leverage that established expression to (i) decompose the current into pseudospinorial components and (ii) quantify cross-correlations induced by the Rashba interaction.
Thus, the role of \eqref{eq:componentes} here is not to introduce a new formula, but to provide a rigorous and practical framework for dissecting the transmission variations and their spinorial origin\cite{Serna2019}, thereby informing future theoretical and experimental studies on spintronic transport in graphene.

Based on these observations regarding the electronic transmission coefficients with and without Rashba interaction, we can synthesize below the main conclusions reached in this work, highlighting their theoretical and technological implications, as well as the perspectives for future studies and applications.
